% ── SECTION 6: THE NON-COMPOSITE SOUL AND DECOHERENCE (~500 words) ──────────

\section{The Non-Composite Soul: Implications}
\label{sec:soul}

% \bahai: soul is not a combination of elements (\AbdulBaha, SAQ).
% What this implies within a relational framework: no internal structure
% to lose coherence between.
% NOT a claim that the soul is a quantum system.
% A structural observation about what non-compositeness entails.

\AbdulBaha teaches that the soul ``is not a combination of elements,
it is not composed of many atoms''~\citep{SAQ}.
The soul is non-composite: it has no constituent parts, no internal structure
from which it could be assembled or into which it could be dissolved.

Within the relational ontological framework established in this paper,
this claim has a precise implication.
Quantum decoherence is a process by which a system's internal degrees of freedom
become entangled with environmental degrees of freedom, causing the system to
lose quantum coherence.
But decoherence requires \emph{internal structure}: there must be parts
between which coherence can be lost.

A genuinely non-composite entity has no internal degrees of freedom.
It therefore cannot undergo quantum decoherence in the technical sense.

This is not a claim that the soul is a quantum system, or that the soul exists
in superposition, or that it can be described by a wave function.
It is a structural observation: if the soul is what the \bahai writings say it is
--- non-composite, not assembled from elements --- then the standard decoherence
argument does not apply to it.

Two implications follow, and both are best stated as questions opened rather
than as positions defended.

The first concerns the long-standing decoherence objection in the quantum
mind debate.
Critics of quantum-mechanical accounts of consciousness have argued that the
warm and noisy environment of the brain destroys quantum coherence too rapidly
for any cognitive process to depend on it~\citep{Zurek2003}.
The objection is decisive against any account that locates the relevant
coherence in macromolecular structures whose decoherence times are short.
It does not, however, apply with the same force to an entity that has no
internal structure at all.
Whether the \bahai conception of the soul can do work in this debate is a
question for further inquiry; what \S\ref{sec:soul} establishes is only that
the standard decoherence objection cannot be redirected against the
non-composite soul without modification, since the objection presupposes
exactly the internal compositeness the \bahai conception denies.

The second concerns the role of the observer.
Relational quantum mechanics makes the relational interaction between system
and observer constitutive of what becomes definite~\citep{Rovelli1996}.
If a non-composite, relationally embedded entity is what \bahai scholarship
has in mind by the soul, then such an entity is structurally available as an
observer in the precise relational sense the physics requires.
The argument does not establish that the soul plays this role.
It establishes only that the structural conditions for it to do so are not
incoherent.
A research programme oriented to this question would draw on \bahai
philosophical anthropology, on the relational reading of quantum mechanics,
and on the conservation arguments developed in \S\ref{sec:convergences}.
We mark the programme as open and leave its substantive investigation to
subsequent work.
